\section{Related Work}
\label{sec:relatedwork}

\subsection{Modeling of BLDC Motor}
The electromechanical model of a BLDC (Brushless DC) motor is foundational for understanding its behaviour under different control schemes. BLDC motors are categorized by their back-electromotive force (back-EMF) waveform: trapezoidal or sinusoidal. This distinction is crucial, as the trapezoidal shape inherently leads to torque ripple when the supplied phase currents are not perfectly aligned, directly influencing the choice and effectiveness of the control strategy \cite{patil_analysis_2025}. For a BLDC motor with trapezoidal back-EMF, the electromagnetic torque is given by:
\[
T_e = \frac{e_a i_a + e_b i_b + e_c i_c}{\omega_m}
\]
where \( e_x \) is the back-EMF and \( i_x \) is the phase current \cite{li_quantitative_2019}. The classical d-q reference frame model, ideal for sinusoidal machines, is less suitable for trapezoidal BLDC motors because it assumes sinusoidal flux distribution. Phase-variable modelling in the natural (abc) frame is therefore more appropriate, as it directly accounts for the non-sinusoidal, trapezoidal nature of the back-EMF and the associated harmonics \cite{mohammd_taher_new_2021}.

\subsection{Trapezoidal Commutation (Six-Step Control) for BLDC Motors}
% Trapezoidal commutation, or six-step control, uses Hall-effect sensors to synchronize phase current switching every 
% 120 electrical degrees. 
Trapezoidal commutation, or Six-Step control, uses bipolar conduction, with two motor phases conducting at any time and current commutation occurring every 120 electrical degrees \cite{gieras_modern_2023}. As commutation depends on rotor position, Six-Step control requires either position sensors (e.g. Hall sensors, encoders, or resolvers) or sensor-less estimation based on back-EMF detection or observers \cite{gieras_modern_2023, gasc_conception_2004}.  This method is renowned for its simplicity of implementation and low hardware cost \cite{bhatiya_bldc_2024}. It enables effective torque control but introduces significant torque ripple during commutation events, especially under high load \cite{jomsa-nga_torque_2024}. This ripple generates noise, increases mechanical stress, and reduces overall efficiency \cite{mohammd_taher_new_2021}. Although PWM techniques can mitigate this ripple, they do not completely eliminate it \cite{li_quantitative_2019}.

\subsection{Field-Oriented Control (FOC) for BLDC Motors}
FOC is a vector control strategy that decouples the stator flux and torque components. It transforms three-phase currents into orthogonal \( I_d \) and \( I_q \) components, enabling precise torque control and significant ripple reduction \cite{jomsa-nga_torque_2024}. FOC is particularly effective for BLDC motors with sinusoidal back-EMF but can also be applied to trapezoidal back-EMF motors, albeit with less impressive ripple suppression results \cite{li_quantitative_2019}. It requires greater computational power and more precise position sensors (e.g. encoders). Comparative analysis shows that FOC yields a more stable stator current profile and significantly reduces torque variations compared to trapezoidal control \cite{patil_analysis_2025}.


\subsection{Comparative Analysis: FOC vs. Trapezoidal for Light Electric Vehicles}

\subsubsection{Torque Ripple and User Comfort}
Firstly, torque ripple can be reduced for both control methods by selecting appropriate motor parameters, such as the number of stator slots and rotor poles \cite{gasc_conception_2004}.
FOC substantially reduces torque ripple compared to Six-Step control, directly enhancing ride comfort and minimizing vibrations. Experimental results show a torque ripple of \SI{18.38}{\percent} for FOC versus \SI{35.67}{\percent} for Six-Step control at 500~rpm \cite{jomsa-nga_torque_2024}. Commutation torque ripple (CTR), prominent in Six-Step control, can be specifically targeted and mitigated using advanced control techniques like Model Predictive Control (MPC) while retaining the fundamental simplicity of trapezoidal commutation \cite{mohammd_taher_new_2021}.

\subsubsection{Energy Efficiency}
FOC optimizes torque per ampere (MTPA), improving efficiency at low loads. Six-Step control exhibits lower switching losses at high speeds \cite{li_quantitative_2019}.

\subsubsection{Complexity, Cost, and Low-Tech Suitability}
Six-Step control is inherently simpler, cheaper, and more robust, making it a prime candidate for low-tech applications. Research focused on reducing propulsion system costs proposes simplified hardware topologies, such as 4-switch inverters (instead of 6) coupled with direct current control strategies, maintaining acceptable performance while significantly lowering hardware costs \cite{lee_advanced_2001}. FOC, while superior in performance, is more complex to implement and carries higher hardware costs (sensors, processing power).

\subsubsection{Dynamic Response}
FOC provides faster response times and better load disturbance rejection \cite{jomsa-nga_torque_2024}.
